\documentclass[a4paper,
               11pt,
               parskip=half,
               headinclude,
               titlepage=false]{scrartcl}

%%%%%%%%%%%%%%%%%%%%%%%%%%%%%%%%%%%%%%%%
\newcommand{\metaauthor}{moPsy}
\newcommand{\metatitle}{MisterClippy User Manual}
\newcommand{\metarev}{1.1}
%%%%%%%%%%%%%%%%%%%%%%%%%%%%%%%%%%%%%%%%


\usepackage[utf8]{inputenc}

\usepackage[T1]{fontenc}        % Tries to use Postscript Type 1 Fonts for better rendering
\usepackage{lmodern}            % Provides the Latin Modern Font which offers more glyphs than the default Computer Modern
\usepackage[sfdefault]{carlito}
\usepackage[intlimits]{amsmath} % Provides all mathematical commands
\usepackage[protrusion=true,
            expansion,
            kerning=true,
            babel=true]{microtype}
\usepackage{fontawesome}

\usepackage[english]{babel}
\usepackage[english]{isodate}

\usepackage{eurosym}
\usepackage[shortlabels]{enumitem}

\usepackage{grffile}            % Allow you to include images (like graphicx). Usage: \includegraphics{path/to/file}

\usepackage[ugly]{units}        % Allows you to type units with correct spacing and font style. Usage: $\unit[100]{m}$ or $\unitfrac[100]{m}{s}$

\usepackage{xspace}             % Use \xpsace in macros to automatically insert space based on context. Usage: \newcommand{\es}{ESPResSo\xspace}
\usepackage[dvipsnames,table]{xcolor}             % Obviously colors. Usage: \color{red} Red text
\usepackage{booktabs}           % Nice rules for tables. Usage \begin{tabular}\toprule ... \midrule ... \bottomrule


%\usepackage[hang]{subfigure}  % unterteilte Abbildungen
\usepackage{setspace}          %fuer dem Zeilanabstand
\usepackage[hang,bf,footnotesize]{caption}   % bessere Bildunterschriften
\usepackage{wrapfig}

\usepackage{csquotes}


\usepackage[hyphens]{url}                % Lets you typeset urls. Usage: \url{http://...}
\usepackage[pdftex,colorlinks=true,linkcolor=black,citecolor=black,urlcolor=black]{hyperref}
\hypersetup{
  pdftitle    = {\metatitle},
  pdfsubject  = {Assembly Guide},
  pdfauthor   = {\metaauthor},
  pdfkeywords = {} ,
  pdfcreator  = {pdflatex},
  pdfproducer = {LaTeX with hyperref}
}
%%%%%%%%%%%%%%%%%%%%%%%%%%%%%%%%%%%%%%%%

\usepackage{pdfpages}

\usepackage{tabularx}


\usepackage[headsepline, footsepline]{scrlayer-scrpage}
\setlength{\headheight}{34pt} % logo height
\clearpairofpagestyles
\ohead{
\vfill%
\begin{tabular}{@{}l l@{}}
Patrick with \includegraphics[height=2em]{moPsy_logo}\hspace*{-0.1cm}
\end{tabular}
}
\ihead{
\vfill%
{\huge MisterClippy} \,v\metarev\, User Manual%
}
\ofoot{Page \thepage}
\setkomafont{pageheadfoot}{\sffamily\footnotesize}
\setkomafont{pagination}{}
\pagestyle{scrheadings}


%\usepackage{background}

%\usepackage[section]{placeins} % keep figures in their section


\definecolor{level_easy}{rgb}{0.2, 0.6, 0.1}

\definecolor{row1}{rgb}{1.0, 1.0, 1.0}
\definecolor{row2}{rgb}{0.93, 0.93, 0.99}



\begin{document}


\begin{minipage}[]{11cm}
\setlength{\parskip}{\medskipamount}
\section*{Thank you for choosing the MisterClippy module!}

This manual provides information on installing and using the MisterClippy module.
We hope that all questions will be answered.

If you encounter a problem or have further questions,
please do not hesitate to reach out to us at \href{mailto:contact@mopsy-music.de}{\texttt{contact@mopsy-music.de}}.


\vspace{1em}
\begin{center}
Enjoy your module! \quad \includegraphics[height=0.8em]{peace_love_music}
\end{center}


\section*{About MisterClippy}

MisterClippy is a clip detector with adjustable threshold and gate output.
And since this is so much fun, we have doubled the feature set in one module!

\end{minipage}%
\hspace{0.5cm}
\begin{minipage}[]{3.5cm}
\includegraphics[width=3.5cm]{MisterClippy-case-shot-frame}
\end{minipage}


\section*{Installation}

The MisterClippy module will occupy \textbf{\unit[12]{HP}} in your rack with a depth of \textbf{\unit[20]{mm}}.

To \textbf{connect power} …
\begin{enumerate}
 \item Please make sure that your cabinet is \textbf{disconnected from the mains} and has no power.
 \item The module is connected with a \textbf{16-pin} (2x8) bus cable.
 \item \textbf{Check your power cable and the polarity.} The coloured ({\color{red}red}) line on the cable (pin 1) needs to be connected to the \textbf{\unit[-12]{V}} rail.
 If the module is installed in an upright orientation, the \textbf{\color{red}red stripe} needs to be \textbf{facing down}.
 \item If you plug the module backwards, it may burn out.
 Unfortunately this cannot be covered by the warranty.
\end{enumerate}

%The module draws a maximum of \textbf{\unit[40]mA} from the \textbf{\unit[+12]{V}} rail.
\emph{Detailed power measurements are pending.}


\begin{minipage}[]{4cm}
\vspace{-2cm}
\includegraphics[width=3.8cm, trim={0 5.5cm 0 2.5cm}, clip]{MisterClippy-calibration}
\end{minipage}
\hspace{0.5cm}
\begin{minipage}[b]{10.5cm}
\setlength{\parskip}{\medskipamount}


\section*{Calibration}

The module has two trim potentiometers, \textbf{\color{red}RV3} and \textbf{\color{red}RV4}, to adjust the default clip detection threshold.
These should be set so that their middle pin presents \textbf{\unit[+10]{V}} with respect to ground.

\end{minipage}


\begin{minipage}[]{11cm}
\setlength{\parskip}{\medskipamount}
\section*{Frontpanel}

The module is divided into two equal and independent sections.

The \textbf{(audio) signal} is fed into the input \textbf{\color{red}IN} (1).

As this signal is usually needed somewhere else, it is \textbf{mirrored} to the output \textbf{\color{red}THRU} (2) without modification.

\vspace{1em}

The \textbf{clip threshold} can be provided as a CV on the \textbf{\color{red}THRSH} input (4) and \textbf{defaults to \unit[+10]{V}}.

In addition, the clip threshold---internal or externally provided---can be \textbf{attenuated} with the potentiometers (3).

\vspace{1em}

Clipping is \textbf{detected} if the signal is \textbf{above} the clip threshold \textbf{or below} the negated clip threshold.
For example, with a clip threshold of \unit[+4]{V}, any voltage greater then \unit[+4]{V} and lower than \unit[-4]{V} will be detected as clipping.

If clipping is detected, the \textbf{\color{red}LED} (5) will light up and a \textbf{gate} is presented at the \textbf{\color{red}CLIP} output (6).

Note that the clip signal will \textbf{persist} for a short time. If another clipping occurs withing this time window, the gate will be \textbf{extended} instead of being re-triggered.

\end{minipage}%
\hspace{0.5cm}
\begin{minipage}[]{3.5cm}
\includegraphics[width=3.5cm]{MisterClippy-frontpanel-labels}
\end{minipage}

\section*{Online Resources}

More about the module  …
\begin{itemize}[noitemsep]
 \item on GitHub: \url{https://github.com/rainerhassel/MisterClippy}
 \item via e-mail: \href{mailto:contact@mopsy-music.de}{\texttt{contact@mopsy-music.de}}
\end{itemize}


\end{document}
